\chapter{Introduction}

\section{Context}
Environmental education is increasingly recognised as an essential component of sustainable development, as it promotes the knowledge, attitudes, and competencies required to understand and address environmental challenges. UNESCO highlights the importance of integrating sustainability and climate action into educational practices, promoting learning experiences that encourage critical thinking and environmental responsibility \cite{unesco_2024}. In this context, direct contact with nature has been shown to support environmentally responsible attitudes and behaviours, particularly when learners engage in active and meaningful experiences \cite{rosa2018nature}.

One example of this approach is the educational activity carried out in the subject MI1815 - The Natural Environment in Early Childhood Education at Jaume I University. The activity takes place in the \textit{Jardín de los Sentidos}, located on the Riu Sec campus and recognized as one of the university’s three main botanical areas of interest. The garden is a botanical and cultural space of approximately 1,500 m², organised into six thematic platforms: a welcome area and five areas related to sensory perception. It contains more than 260 plant species belonging to 78 different families. Established in 2004 as both an educational and cultural resource, the garden serves as an outdoor laboratory where future teachers can engage directly with the natural environment.

Through guided exploration of the garden, students observe a wide variety of plant species from different regions of the world, develop an appreciation for botanical heritage, and learn to interpret ecological relationships within a real-world context. The sensory-oriented design of the garden also provides opportunities to connect botanical knowledge with pedagogical practices relevant to early childhood education. In addition, the activity encourages students to design and reflect on educational itineraries in natural settings, highlighting the value of nearby green spaces as accessible learning environments that can support environmental education when access to larger natural areas is limited.


\section{Motivation}
Although the \textit{Jardín de los Sentidos} activity provides valuable educational experiences, there are currently several limitations to its impact. As it takes place outdoors and within a specific academic timeframe, its implementation is constrained by the duration of the session, environmental conditions, group availability, and the presence of an expert instructor capable of guiding observations, answering questions, and adapting the experience to the learners’ needs. These constraints make it difficult to repeat the activity frequently, extend it to additional groups, or maintain it as an ongoing learning process. Consequently, there is a need to explore alternative approaches that preserve the experiential value of the activity while enabling more autonomous, flexible, and sustained participation.

In this context, serious games and mobile technologies offer a promising approach for extending such learning experiences into more accessible, interactive, and engaging educational formats. Serious games are digital games designed with a main purpose that goes beyond entertainment. Although they preserve essential elements of play, such as rules, challenges, interaction and feedback, their design is oriented towards educational, training, social or behavioural objectives \cite{hands2026serious, facchino2025seriousgames}.

In educational contexts, serious games are especially relevant because they can place learners in an active role. Instead of receiving information passively, players interact with content, make decisions, test hypotheses and progressively construct meaning through experience. This approach is closely related to constructivist and experiential learning perspectives, as the game environment can encourage exploration, problem solving and reflection. In addition, serious games can offer safe spaces in which learners are allowed to make mistakes and observe their consequences without the risks associated with real-world practice \cite{gomez2019mixed, hands2026serious, facchino2025seriousgames}.

Another key contribution of serious games is their capacity to provide immediate feedback. The system can respond directly to the player’s actions, allowing learners to adjust their strategies, correct errors and continue progressing through an iterative process. This feedback can also support assessment, not as an external interruption of the activity, but as part of the game experience itself. However, the educational value of serious games depends on their design and integration within a broader learning context. Poor interface design, excessive cognitive load, technical problems or lack of pedagogical alignment can reduce their effectiveness. Therefore, serious games should be understood as complementary tools rather than replacements for traditional teaching methods \cite{de2014serious, facchino2025seriousgames, hands2026serious}.

The increasing interest in serious games is also linked to technological development. Recent advances in mobile devices, mixed reality, augmented reality, virtual reality, artificial intelligence and learning analytics have expanded the possibilities for creating interactive educational experiences situated in physical or simulated environments. In this context, mobile and location-based games are particularly relevant because they can connect digital interaction with real-world exploration, making the environment itself part of the learning process \cite{gomez2019mixed, hands2026serious}.

By combining these approaches, guided outdoor educational activities can be transformed into autonomous and interactive learning experiences in which students engage with challenges, receive feedback, collaborate with peers, and construct environmental knowledge through the exploration of real-world environments.

From this perspective, the present project emerged from a collaboration between the Department of Early Childhood Education and the CEVI research group, to which the supervisor of this Bachelor’s Thesis belongs. The aim is to transfer the core learning outcomes of the original activity into a mobile experience that preserves interaction with the real environment while reducing the need for continuous teacher guidance and the constraints associated with scheduled learning sessions. By integrating geolocation, collaborative gameplay and situated learning principles, the proposed solution (\textit{Micorrizas}) aims to provide a more autonomous, scalable and engaging educational experience that complements the original activity and supports environmental education beyond the classroom.

This project is also motivated by professional and personal development. From a technical perspective, it represents a challenge because it covers areas I had not previously developed in my university projects, particularly multiplayer interaction and geolocation-based gameplay. The implementation of a serious mobile game that combines map integration, location-based triggers, cooperative synchronisation and mini-games offers the opportunity to expand my skills as a video game developer. At the same time, the project is personally motivating because it connects game design with a real educational context, allowing the final prototype to address a specific need and explore how video games can support learning.

\section{Objectives}

The main objective of this project was to design and develop an educational video game that promotes environmental awareness and engagement through active exploration and collaborative learning.

The specific objectives to achieve the main goal are the following:
\begin{enumerate}
    \item To transform the educational activity conducted in the \textit{Jardín de los Sentidos} into a videogame while preserving its pedagogical contents.
    \item To design cooperative multiplayer gameplay mechanics that encourage exploration, observation, communication, and collaborative problem-solving.
    \item To develop an Android mobile game that reacts to the player's physical position within the \textit{Jardín de los Sentidos}.
    \item To provide teachers and the students of the subject an educational resource that complements traditional guided visits and promotes autonomous learning.
    \item To increase awareness and appreciation of biodiversity through interactions within the natural environment.
\end{enumerate}

